\documentclass[10pt,a4paper]{amsart}
\usepackage[margin=19mm]{geometry}
\usepackage{amsmath,amssymb,mathtools}
\usepackage[scr=rsfs,cal=euler]{mathalfa}
\usepackage{xfrac}
\usepackage[unicode,hidelinks,bookmarks=false]{hyperref}
\newcommand{\ie}{i.e.\ }
\newcommand{\cf}{cf.\ }
\newcommand{\resp}{respectively\ }
\newcommand{\Subsection}{\subsection}
\providecommand{\nobreakdash}{\nobreak\hbox{-}}
\setlength{\emergencystretch}{2em}
\multlinegap0pt
\usepackage{bm}
\usepackage{bbm}
\usepackage{dsfont}
\usepackage{xspace}
\usepackage{cancel}
\usepackage{mathtools}
\usepackage{amsthm}
\makeatletter
\@ifundefined{theorem}{\newtheorem{theorem}{Theorem}}{}
\@ifundefined{proposition}{\newtheorem{proposition}[theorem]{Proposition}}{}
\makeatother
\title[Efficient points in a sum of sets of alternatives]{Efficient points in a sum of sets of alternatives}
\author{Anas Mifrani}
\date{Source: arXiv:2412.16331v2 (CC BY 4.0)}
\begin{document}
\maketitle

\noindent\emph{Source and licence.} Efficient points in a sum of sets of alternatives. Version arXiv:2412.16331v2 (CC BY 4.0), distributed under the Creative Commons Attribution 4.0 International licence, as recorded in the arXiv licence metadata for this exact version and shown on the abstract page https://arxiv.org/abs/2412.16331v2. Full rights notice: licenses/arxiv-2412.16331-CC-BY-4.0.txt.\par\smallskip

\noindent\textit{Editorial scope.} Counted unit: the Proposition whose source label is \texttt{prop4} in the article (source0.tex lines 890--907, source label \texttt{prop4}), delivered together with the article's own complete original proof of it (source0.tex lines 909--988, one \texttt{proof} environment of 80 lines). No mathematics is written, repaired or re-derived: statement and proof are the article's own text line for line. Required context retained uncounted: none. Every definition, display and label of those lines is retained unchanged. Omitted from this package and not redistributed: 1--889, 908--908, 989--1511. Nothing inside a retained excerpt is omitted or altered: every retained body is delivered exactly as the article's lines are written, so no omission receipt of any kind accompanies this package. Citations: the retained excerpts carry no citation command at all, so no bibliography entry of the article is transcribed and no reference is redistributed. Reference adaptation: none was needed and none was made; every reference inside the delivered unit resolves inside it, so no unresolved reference or citation is produced. Printed slips kept literal: no printed word, symbol, hypothesis or number of the counted statement or of its proof was altered. Layout and class: the article's own class is \texttt{sn-jnl} with packages including enumitem, graphicx, multirow, amsmath, amssymb, amsfonts, amsthm, mathrsfs, appendix, xcolor, textcomp, manyfoot, booktabs, algorithm; the wrapper is the lane's pinned \texttt{amsart} editorial wrapper at 10pt on A4, which restates the theorem-like environments the retained text needs and adds the article's own macro definitions that the retained text needs (none), with extra wrapper packages (bm, bbm, dsfont, xspace, cancel, mathtools). The delivered numbering is the unit's own. Assets: no retained span contains a figure environment, a graphics-inclusion command, a PDF-page inclusion, a table or a TikZ picture; no image file is copied into this package, the package declares no asset dependency and no image file is redistributed with it. Licence: version arXiv:2412.16331v2 of this article is distributed under the Creative Commons Attribution 4.0 International licence, as recorded in the arXiv licence metadata for this exact version and shown on the abstract page https://arxiv.org/abs/2412.16331v2. A full rights notice is delivered at licenses/arxiv-2412.16331-CC-BY-4.0.txt. Characters: every retained excerpt is pure ASCII, so the delivered engine, XeLaTeX with its default Latin Modern text font, reports no missing character.
\begin{proposition}
\label{prop4}
%Suppose, in addition to the hypotheses already imposed on $R$, that $0_{G}P(pb)$ implies $0_{G}Pb$ for all $p \geqq 1$ and $b \in G$. 
Assume (P1), (P3) and (P5). Let $a^1, ..., a^n$ be $n \geqq 2$ distinct points in $G$, and let $b \in B$. Form a set of indices $i_1, ..., i_n$ from $\{1, ..., n\}$ with $i_j \neq j$ for each $j$. Then the system
\begin{equation}
\tag{$S^{4}$}
\label{$S^{4}$}
  \left\{
    \begin{aligned}
      & a^{i_1}P(a^{1} + b), \\
      & \dots \\
      & a^{i_{n}}P(a^{n} + b), \\
      & bI{0_G},
    \end{aligned}
  \right.
\end{equation}
is inconsistent.
\end{proposition}

\begin{proof}
Suppose system (\ref{$S^{4}$}) were consistent. Define the permutation $f:\{1, \dots, n\} \rightarrow \{1, \dots, n\}$ that maps each $j = 1, \dots, n$ to the index $i_j$ which appears on the left-hand side of the comparison $a^{i_j}P(a^{j}+b)$. Fix any index $j_0$ and consider the sequence \begin{equation*}j_0, j_1 := f(j_0), j_2 := f^2(j_0), \dots, j_n:= f^{n}(j_0).\end{equation*}

This sequence says that \begin{equation*}
  \left\{
    \begin{aligned}
      & a^{j_{1}}P(a^{j_0} + b), \\
      & a^{j_{2}}P(a^{j_{1}}+b), \\
      & \dots \\
      & a^{j_n}P(a^{j_{n-1}} + b), \\
    \end{aligned}
  \right.
\end{equation*} 

Now each element in the sequence ranges in $\{1, \dots, n\}$, while the sequence itself runs to $(n+1)$ members. By the pigeonhole principle, therefore, there must exist a pair of integers $0 \leqq r < m$ with $j_r = j_m$. The subsequence starting at $j_r$ and culminating in $j_m$ forms a cycle. Indeed, because this subsequence means \begin{equation*}
  \left\{
    \begin{aligned}
      & a^{j_{r+1}}P(a^{j_r} + b), \\
      & a^{j_{r+2}}P(a^{j_{r+1}}+b), \\
      & \dots \\
      & a^{j_m}P(a^{j_{m-1}} + b), \\
    \end{aligned}
  \right.
\end{equation*}
holds, and inasmuch as $j_m = j_r$, iterative application of transitivity (P1) and isotonicity (P3) to these comparisons (beginning with the bottom comparison and ascending) generates the comparison \begin{equation*}a^{j_r}P(a^{j_r}+(m-r+1)b).\end{equation*} When we apply property (P3) to this very relation, we find that \begin{equation*}0_GP(m-r+1)b,\end{equation*}
so that $0_GPb$ by property (P5). But this cannot be so, since $bI0_{G}$. We infer from this that (\ref{$S^{4}$}) is inconsistent, precisely as the proposition avers. 
%We proceed by induction on $n$. For $n = 2$, the only candidates for $i_1$ and $i_2$ are $i_1 = 2$ and $i_2 = 1$. Suppose, contrary to the proposition, that we had simultaneously $a^{2}P(a^{1} + b)$, $a^{1}P(a^{2}+b)$ and $bI0_{G}$. We have assumed property (P3). %Thus we have
%$a^2P(a^1+b)$ and $a^1P(a^2+b)$.

%It would therefore follow that $((-a^1)+a^{2})Pb$ and $((-a^{2})+a^1)Pb$, and this would imply $0_{G}P(2b)$, and hence $0_{G}Pb$ by property (P5). That would be absurd in light of $bI0_{G}$.


%From the first relation we derive
%$((-a^1) +a^2)Pb$, and from the second $((-a^2) + a^1)Pb$. Noting that the group inverse of $(-a^1) + a^2$ is $(-a^2) + a^1$, direct application of (P3) on the latter of the relations just obtained yields the comparison $0_GP((-a^1)+a^2 + b)$. Another application of (P3), this time on $((-a^1)+a^2)Pb$, results in $((-a^1) + a^2 + b)P(2b)$. Since $R$ is transitive (P1), the previous two sentences lead to the conclusion that $0_GP(2b)$; that is, $0_GPb$, as per property (P5). We know, however, that this cannot be true, since $bI0_G$.  %Now translate both comparisons (again using (P3)) so that their left-hand sides align: from 
%$(a^2+(-a^1))Pb$ we get 
%\[
%0_GP(a^1+(-a^2) + b),
%\]
%and from $(a^1+(-a^2))P b$ we get 
%\[
%(a^1+(-a^2) + b)P(a^1+(-a^2) + 2b).
%\]
%The system is therefore inconsistent in the case $n=2$.


%For the inductive step, let $n \geqq 2$ be an integer such that for any $n$ points $a^1, ..., a^n$ in $G$ and each set of indices $i_1, ..., i_n$ satisfying $i_j \neq j$, we do not have
%\begin{equation*}
 % \left\{
  %  \begin{aligned}
   %   & a^{i_1}P(a^{1} + b), \\
    %  & \dots \\
    %  & a^{i_{n}}P(a^{n} + b), \\
     % & bI{0_G},
    %\end{aligned}
  %\right.
%\end{equation*}
%for any $b \in G$. Then it is clear that for any $(n+1)$ points $a^1, ..., a^{n+1} \in G$, any choice of indices $i_1, ..., i_{n+1}$ satisfying $i_j \neq j$ for each $j = 1, ..., n+1$, and for any $b \in G$, the system
%\begin{equation*}
 % \left\{
  %  \begin{aligned}
   %   & a^{i_1}P(a^{1} + b), \\
    %  & \dots \\
     % & a^{i_{n+1}}P(a^{n+1} + b), \\
      %& bI{0_G},
    %\end{aligned}
  %\right.
%\end{equation*}
%is inconsistent because it implies the system
%\begin{equation*}
 % \left\{
  %  \begin{aligned}
   %   & a^{i_1}P(a^{1} + b), \\
    %  & \dots \\
     % & a^{i_{n}}P(a^{n} + b), \\
    % & bI{0_G},
    %\end{aligned}
  %\right.
%\end{equation*}
%which we know is inconsistent on account of the induction hypothesis. The proof is complete.
\end{proof}

\end{document}
